\section{Problem Formulation}\label{sec:formulation}

% Three-policy-lane view of P-CONSENT. The three policies are the primary
% structure: pi^DR (day-ahead), pi^ch (every dt), pi^neg (per arrival), each
% on its own clock. Vertical arrows are the within-day dependency (FSL ->
% charging feasible set -> negotiation options); the dashed arrow is the sole
% cross-day link (surviving buffer). Not a timeline: see Sec. 4.
% Requires \usetikzlibrary{positioning, fit, backgrounds, arrows.meta}.
\begin{figure*}[t]
\centering
\newcommand{\ltitle}{\bfseries}%
\newcommand{\leqn}{\scriptsize\itshape\color{black!55}}%
\begin{tikzpicture}[
  font=\footnotesize,
  every node/.style={transform shape},
  lane/.style={rectangle, rounded corners=3pt, draw=black!55, thick,
               align=left, inner sep=7pt, text width=0.60\textwidth,
               minimum height=1.55cm},
  clock/.style={rectangle, rounded corners=3pt, draw=black!40, thick,
                fill=white, align=center, inner sep=4pt, text width=1.9cm,
                minimum height=1.55cm, font=\scriptsize},
  dep/.style={-{Latex[length=2.2mm]}, very thick, black!75},
  back/.style={-{Latex[length=2.2mm]}, very thick, dashed, black!75},
  deplab/.style={font=\scriptsize\itshape, text=black!60, align=center, fill=white, inner sep=1pt}
]
% ============ LANE 1: pi^DR (day-ahead) ============
\node[clock, fill=blue!8] (cDR) {\ltitle Day-ahead\\(end of $d{-}1$)\\[2pt]\leqn once/day};
\node[lane, fill=blue!5, right=0.35cm of cDR] (DR) {%
{\ltitle Layer $\pidr$: demand-response commitment.}\ \ Forecast the surviving
buffer, then commit a defensible peak cap $\FSL_d$ and earn a capacity payment.
Fixed before any day-$d$ arrival.\\[2pt]
{\leqn $\psihat\!\to\!\buffhat{d}\!\to\!\Omeps$; set
$\FSL_d(\lambda)$;\ earn $\Cinc{d}$}};

% ============ LANE 2: pi^ch (every dt) ============
\node[clock, fill=teal!9, below=1.7cm of cDR] (cCH) {\ltitle Every $\dt$\\[2pt]\leqn continuous};
\node[lane, fill=teal!6, right=0.35cm of cCH] (CH) {%
{\ltitle Layer $\pich$: real-time charging.}\ \ Set charge/discharge rates at
least cost under the committed FSL; bank only into spare capacity; during an
event (if called) discharge to hold load under $\FSL_d$.\\[2pt]
{\leqn minimize s.t.\--;
bank if $\Ptot{t}\!<\!\min(\Pmaxhat,\text{cap})$; slack $\vio{t}$ via $\Cpen{d}$}};

% ============ LANE 3: pi^neg (per arrival) ============
\node[clock, fill=orange!10, below=1.7cm of cCH] (cNEG) {\ltitle Per arrival\\/ departure\\[2pt]\leqn event-driven};
\node[lane, fill=orange!6, right=0.35cm of cNEG] (NEG) {%
{\ltitle Layer $\pineg$: negotiation \& persistence.}\ \ Offer each driver
compensated flexibility and a banking price $\pbank{v}{k}$; at departure settle
energy, banking, and realized event/peak help; carry the surviving buffer over.\\[2pt]
{\leqn price by $\Uflex$$\Rightarrow\!\Eneg,\tdep$; settle
$\ucost{k}{v}$, $\Ubuf,\Udc$--;
carry $\survived{v}{k+1}$}};

% ============ WITHIN-DAY DEPENDENCY (top -> down) ============
\draw[dep] (DR.south) -- (CH.north)
  node[deplab, midway] {$\FSL_d$ constrains charging feasible set};
\draw[dep, <->] (CH.south) -- (NEG.north)
  node[deplab, midway] {discharge capability sets offerable flexibility};

% ============ CROSS-DAY LINK (bottom -> up, dashed) ============
\draw[back] (NEG.east) -- ++(0.7,0)
  node[midway, right, font=\scriptsize\itshape, text=black!70, align=center]
    {surviving buffer\\ $\buff{d}\!\to\!\buffhat{d+1}$}
  |- (DR.east);

  

\end{tikzpicture}
\caption{The three policies of P-CONSENT, each acting on its own clock:
$\pidr$ commits the firm service level a day ahead, $\pich$ sets charging and
discharging rates every $\dt$, and $\pineg$ negotiates flexibility and settles
payments at each arrival and departure. Solid vertical arrows are the sole
within-day dependency, the committed FSL constrains the charging feasible set,
which in turn sets the flexibility the negotiation can offer; this is a
dependency, not a temporal order, since the charging and negotiation epochs
interleave continuously through the day (the problem definition (Section~\ref{sec:formulation})). The dashed
arrow is the only coupling between days: the buffer that survives a driver's
trip becomes tomorrow's deliverable buffer and sharpens the forecast that sets
the next commitment.}
\label{fig:processflow}
\end{figure*}


We consider a single commercial building that manages its own EV charging infrastructure, serves a population of persistent EV users, and participates in a utility demand response (DR) program. Three decisions act on the system at three timescales (Figure~\ref{fig:processflow}): a day-ahead \emph{commitment} of a Firm Service Level to the utility (policy $\pidr$), real-time \emph{charging} and discharging (policy $\pich$), and per-arrival \emph{negotiation} of compensated adjustments to a user's required SoC and departure time (policy $\pineg$). We first declare the shared physical layer (Section~\ref{sec:specs}), then define each of the three sub-problems, their decisions, constraints, and objective contribution. 
% Section~\ref{sec:coupled} then assembles them into the single sequential problem they jointly approximate and states its ideal solution. 
The complete variable dependency structure is summarized in Table~\ref{tab:vartree}.








\subsection{Specifications}\label{sec:specs}

\paragraph{Time structure.} The billing cycle consists of days $d \in \Dset$. Each day is discretized into time steps $t \in \Td$ of length $\dt$ (e.g., 15 minutes), and $\Tset = \cup_d \Td$. The demand charge is assessed over designated peak periods $\Tpeak \subseteq \Tset$. Control decisions are taken at the start of each step and held constant for its duration.

\paragraph{Persistent users and sessions.} Let $\Vset$ be the set of EV
users, each with a consistent identity across the billing cycle. User $v$'s visits are indexed by sessions $k \in \Kv$, with first session $k = 0$; session $k$ spans $[\tarr(v,k), \tdep(v,k)] \subseteq \Tset$, and $\Tvk$ denotes its time steps. Sessions need not occur on consecutive days; the formulation accommodates arbitrary gaps between visits. At plug-in, the arrival time $\tarr(v,k)$ and arrival SoC $\Earr(v,k)$ are observed, and the user reports a required departure SoC $\Ereq(v,k)$ and departure time $\treq(v,k)$. Let $\Vd{d} \subseteq \Vset$ denote the users with a session active on day $d$. Per-user, per-session quantities are written $X(v,k)$ (and $X(v)$ when they depend on the user alone); within a report or agreement tuple that concerns a single fixed user-session, we drop the arguments and write, for instance, $(\Ereq, \treq)$ for brevity.

\paragraph{Battery state and persistence.} The SoC of EV $v$ at time $t$
is $\soc{t}{v}$, expressed in energy units (kWh) and bounded by the usable battery window:
\begin{equation}\label{eq:socbounds}
\Emin(v) \;\le\; \soc{t}{v} \;\le\; \Emax(v).
\end{equation}
Within a session, the SoC evolves as
\begin{equation}\label{eq:soc}
\soc{t+1}{v} \;=\; \soc{t}{v} + \rate{t}{v}\,\dt,
\end{equation}
where $\rate{t}{v}$ (kW) is the charging rate, negative for discharge; since at most one session is active at any time, the session index is omitted from per-step quantities. The EV user arrives with a SoC requirement for the session $\Ereq(v,k)$ (kWh). As part of negotiations, the users can accept a different negotiated SoC $\Eneg(v,k)$. This is separate from the SoC at departure $\Edep(v,k)$ (kWh). Between sessions $k$ and $k{+}1$, the EV consumes $\Eext(v,k)$ energy off-site (kWh), net of any external charging, with $\Eext(v,k) \le \Edep(v,k)$. Combined, 
\begin{equation}\label{eq:carryover}
\Earr(v,k{+}1) \;=\; \Edep(v,k) - \Eext(v,k),
\end{equation}
where $\Edep(v,k)$ is the SoC at departure from session $k$. The arrival SoC is therefore exogenous only at $k = 0$; for every later session it is determined by \eqref{eq:carryover}, which is precisely what makes the fleet an inter-temporally coupled resource and what single-day formulations cannot represent. 

\begin{assumption}[Exogenous off-site usage]\label{as:exo}
$\Eext(v,k)$ is unaffected by the over-charge decision; cost-free energy does not induce additional vehicle usage. This is a behavioral assumption subject to a potential rebound effect, revisited when interpreting the measured yield (Section~\ref{sec:eval}).
\end{assumption}


\paragraph{Chargers and assignment.} The building operates a set
$\mathcal{C}$ of heterogeneous charger types with $N_c$ units of type $c$, rate limits $[\rmin{c}, \rmax{c}]$ ($\rmin{c} < 0$ only for bidirectional chargers). Charging is modeled as lossless; discharge delivers energy to the building at round-trip efficiency $\eta \in (0,1]$, so that banked excess of $E$ kWh yields $\eta E$ kWh at the meter. An assignment map $\assign_t : \mathcal{V}_t \leftrightarrow \mathcal{C}_t$ matches present EVs to chargers one-to-one, first-come, first-served, fixed for the duration of the session; an EV that arrives when no charger is free cannot participate. The availability indicator is $\avail{t}{v} \in \{0,1\}$. Charging-rate slowdown at high SoC is modeled by piecewise-linear segments and deferred to the appendix.




\paragraph{Building power and electricity cost.} The building's total
power draw is
\begin{equation}\label{eq:load}
\Ptot{t} \;=\; \bload{t} + \sum_{v \in \Vset} \avail{t}{v}\,\rate{t}{v},
\qquad \Ptot{t} \ge 0,
\end{equation}
where $\bload{t}$ (kW) is the non-EV building load. The energy cost at step $t$ is $\Phi^{e}_t = \we{t}\,\Ptot{t}\,\dt$ with time-of-use price $\we{t}$ (\$/kWh), and the energy cost attributed to user $v$ sums the energy $\gch{t}{v} \ge 0$ delivered to their EV; discharged energy carries no meter credit, its value reaching the user exclusively through the replay incentives of Section~\ref{sec:negotiation}. The demand cost is
\begin{equation}\label{eq:peak}
\Phi^{d} = \wdr{}\cdot \max_{t \in \Tpeak} \Ptot{t}
\end{equation}
Discharging accrues a degradation cost $\Kbatt$ per kWh of throughput, tracked by the one-sided variable $\qdis{t}{v} \ge -\rate{t}{v}\dt$ (nonzero only on discharge, where $\rate{t}{v}<0$; pinned to zero on charge by the objective), and departing below the negotiated floor incurs $\Ksoc$ per kWh of shortfall $\zund{k}{v}$. The building can build towards and forecast a realizable SoC buffer $\buffhat{d}$ for the upcoming DR day, using the EVs present in the current day and that have been present in the past.


\paragraph{Policies.} The building's policy is the triple $\pi =
(\pidr, \pich, \pineg)$ with decision epochs at the end of each day, at every time step, and at every user arrival, respectively. All derived quantities below (SoC trajectories, peaks, costs) are implicitly determined by $\pi$ and the realized user choices; we omit this dependence from the notation for readability.

% \subsection{The FSL Commitment Sub-Problem ($\Pdr$)}\label{sec:dr}

\paragraph{Demand Response Commitment.} For a DR event announced for day $d$, the peak level commitment, Firm Service Level $\FSL_d$ (kW) has to be made ahead of time on day $d{-}n$, where $n$ is the notification time from the power utility. \textcolor{red}{check} The operator is expected to have a tolerance of $\varepsilon$ towards any deviation if the committed FSL. No other decision is taken at this layer. There can be multiple DR events per month, and each has its own commitment and penalties.

\paragraph{Payment and penalty.} The utility registers a baseline
$\Phist$ (kW) at enrollment from historical billing data. The capacity payment and the violation penalty are
\begin{equation}\label{eq:dr}
\Cinc{d} = \winc\,(\Phist - \FSL_d), \qquad
\end{equation}
\begin{equation}
\Cpen{d} = \Kfsl \sum_{t \in \evwin{d}} (P_{t} - \FSL_d) \Delta t \le \sigma^i_t\ 
% ; \quad \forall t \in {\evwin{d}}
%  = \Kfsl \sum_{t \in \evwin{d}} \vio{t}\,\dt,
\end{equation}
where $\winc$ (\$/kW) is the per-day-normalized capacity rate and $\vio{t} = \max(0,\, \Ptot{t} - \FSL_d)$ is the FSL violation, i.e., sum of energy (kWh) used above the FSL during the DR event $\evwin{d}$. The event window $\evwin{d}$ and the demand-charge peak period $\Tpeak$ are distinct time sets and need not coincide.

\begin{remark}[DR Program choice]\label{rem:program}
Deployed programs differ in commitment cadence: the Base Interruptible Program fixes the FSL at enrollment for a season, whereas capacity bidding programs admit nomination with day-ahead and day-of options. We model a day-ahead re-committable FSL, corresponding most closely to the day-ahead capacity-bidding option, and normalize the capacity rate to $\winc$ per committed day. An enrollment-fixed FSL is the constrained special case $\FSL_d \equiv \FSL$ and is reported in the sensitivity analysis.
\end{remark}


\paragraph{User Cost, Choices and options.} User $v$'s choice $\theta_v \in \boldsymbol{\theta}$ records the observed arrival data and the negotiated departure SoC and time. Users are requested to be flexible in their requests and are parameterized by levels $l \in \{0, \ldots, L\}$ with deviation bounds $\Delta e^{\max}_l$ and $\Delta t^{\max}_{\mathrm{dep},l}$; level $l = 0$ is the user's exact request, together with a reject-all choice priced at the external rate $\wext$. Each user has to pay a charging cost $\zeta_v$ depending on the option they choose. 

% \paragraph{Over-charge offer, and the banking price.}
% Independently of the chosen level, the user may consent to over-charge above the negotiated floor, recorded by $\accept{v}{k} \in \{0,1\}$. Consenting is the user's only SoC banking decision, with the amount and timing being controlled by the operator (building), which removes any reportable banking quantity from the user's strategy space. The departing SoC excess $\banked{v}{k}$ is priced at the banking rate
% \begin{equation}\label{eq:pbank}
% \pbank{v}{k} \;=\;  
% \begin{cases}  
% {w}^{e}_{k}, & \text{if DR event not announced} \\  
% 0, & \text{if DR event is announced}  
% \end{cases} 
% \end{equation}
% where ${w}^{e}_{k}$ is the energy price paid.

\paragraph{Total Cost.}  
The building’s total cost is over a monthly billing period under demand response, negotiation, and charging policies, and the user choices. For convenience, we omit explicit notation of this dependence, writing the cost function as
\begin{equation}\label{eq:obj}
% \small
\begin{split}
& J = \sum_{t \in \mathcal{T}} \Phi^{e}_t + \Phi^d
+ K_{\mathbf{SOC}} \sum_{v \in \mathcal{V}} \left(e^v_{\mathrm{req}} - e^v_{\mathrm{dep}} \right)\\
& + K_{\mathbf{batt}} \sum_{v \in \mathcal{V}} \sum_{t \in \mathcal{T}^v} q_t^v + \sum_{d \in \mathcal D} \Cpen{d} - \Cinc{d} 
\end{split}
\end{equation}

subject to the constraints \eqref{eq:socbounds}--\eqref{eq:dr}.


